\newpage
\section*{引言}
高维数据在科学与工程的众多领域中自然出现，涵盖机器学习、信号处理、计算物理与统计学。这类数据常以张量表示，亦即矩阵向高维的推广。张量虽为多模态结构提供了自然的表示，但随着张量阶数增大，直接操作它很快变得困难：参数数目呈指数增长，涉及大量指标的代数表达式也难以理解与实现。

张量网络（tensor network, TN）为应对这些挑战提供了有效的框架。
由~\cite{penrose1971applications} 最早提出、又在量子物理界得到充分发展的张量网络图形语言，把张量收缩编码为图中的边；这一形式体系不仅减少了记号负担，还揭示出被指标记法掩盖的张量运算结构性质。尽管高维张量在现代机器学习与数值分析中居于核心地位，张量网络图在量子计算之外仍未得到充分利用，部分原因是缺少一部自成体系、面向广大技术读者的数学参考书。
本文稿的目标是提供一部关于张量网络及其在张量代数中用法的自洽指南。我们以图示记法讲述张量的主要运算，如收缩、乘积与重排，并说明经典张量分解与相关计算如何在这一框架下自然表达。此外，我们展示张量网络如何简化梯度的推导以及高维概率分布的处理。

贯穿全文，我们说明图示方法并不只是简化记号；它对经典恒等式、秩的界与梯度公式给出了真正更短、更清晰的证明，而这些结果若改用指标运算则需繁复的推导。 



\subsection*{结构安排}

本文稿共分五章，各章以前一章为基础，逐步建立起 
关于张量网络的完整而自洽的论述。

\medskip
\noindent\textbf{第 1 章：张量网络基础。}
我们从第一性原理出发引入张量网络的图形语言。张量表示为带标注腿
（leg）的节点，收缩则编码为共享的边。我们
在这一框架内建立核心运算：内积、外积、迹与 Frobenius 范数，并说明熟知的矩阵恒等式如何自然推广到
高阶张量。本章以复制张量（超边）与矩阵分解的图形
表示收尾，包括 QR 分解与奇异值
分解（SVD）。

\medskip
\noindent\textbf{第 2 章：张量上的运算。}
我们以形式化与图形化两种方式讨论张量的重排运算、纤维、切片、矩阵化与向量化。随后建立主要的张量乘积：模 $n$ 
乘积、Kronecker 积、Khatri-Rao 积与 Hadamard 积，每种乘积都配有其关键恒等式的图示证明。本章以一个一般性定理收尾：借助张量网络图中的割为张量矩阵化的秩给出上界，从而把若干经典秩不等式统一于同一框架之下。

\medskip
\noindent\textbf{第 3 章：张量分解。}
我们介绍高维张量的三种主要分解模型。CANDECOMP/PARAFAC（CP）分解把张量写成若干秩一项之和；Tucker 分解借助压缩的核心张量与因子矩阵将此推广，
而张量列（TT）分解给出链式结构的因子分解，其参数个数随张量阶线性增长。对每种模型，我们给出数学表述、相应的张量网络图、秩的刻画与计算算法，包括用于 Tucker 分解的高阶 SVD（HOSVD），以及用于 TT 的 TT-SVD 与交替最小二乘（ALS）算法。我们进一步讨论 TT 格式下的高效代数运算，并简要概述更先进的架构，包括张量环分解、投影纠缠对态（PEPS）分解与层次 Tucker 分解。

\medskip
\noindent\textbf{第 4 章：计算张量网络的梯度。}
我们为张量网络函数的求导建立一套图示演算。从梯度与雅可比矩阵的经典定义出发，我们证明：张量网络关于任一核心张量的雅可比矩阵，
只需把该核心从图中删去即得。由此得到对经典求导恒等式的透明、纯图示层面的证明，
包括二次型与迹泛函的梯度，并自然推广到同一核心出现多次的网络。我们以两个应用示范这一框架：CP 分解的
基于梯度的优化，以及经由以矩阵乘积算符（MPO）格式参数化的
权重矩阵的反向传播。

\medskip
\noindent\textbf{第 5 章：概率分布与随机张量网络。}
我们探讨张量网络在概率与统计中的
两个相互关联的应用。首先，我们说明离散随机变量的
联合概率分布如何编码为张量，以及边缘分布与条件分布如何化为简单的
图示运算。随后我们讨论张量列格式如何为高维分布提供
高效而结构化的参数化，其中包括源自量子物理的
Born 机器式表示。其次，我们计算元素取自标准正态
分布的随机张量网络的期望。借助 Isserlis 定理与图示推理，我们为
形如
$\EE[\|\Ab\Bb\|_F^2]$ 与
$\EE[(\Ab^\ts \Ab)^{\kron 2}]$ 这类量导出闭式表达式，其简洁程度是常规基于指标的方法难以达到的。

\bigskip
\subsection*{读者对象}

本文稿可作为一部自成体系、面向研究者的参考书，适合
处理高维数据、希望熟练使用张量网络的研究人员与
从业者。所述内容易于入门：只要具备本科第一门线性代数课程层面的
实用知识，包括矩阵乘法、QR 分解与 SVD，
就足以阅读正文的大部分内容。

行文的组织方式同时服务于多个群体。对
\textit{数学家}而言，本文稿给出严格的定义、秩的刻画与
以线性代数理论为基础的算法分析。对
\textit{机器学习与数据科学的研究者}而言，它提供一套实用的
工具，用以参数化高维权重矩阵、经由结构化的张量分解计算梯度，
并以紧凑的形式建模复杂分布。对
\textit{物理学家与量子计算研究者}而言，它把张量网络的图示
约定与应用数学的语言沟通起来。 

本文稿中的所有张量网络图均以 \LaTeX\ 配合 \texttt{TikZ} 宏
排版，图示记法逐步引入，因此不假定读者
此前接触过张量网络图。


